\documentclass[10pt,a4paper]{amsart}
\usepackage[margin=19mm]{geometry}
\usepackage{amsmath,amssymb,mathtools}
\usepackage[scr=rsfs,cal=euler]{mathalfa}
\usepackage{xfrac}
\usepackage[unicode,hidelinks,bookmarks=false]{hyperref}
\newcommand{\ie}{i.e.\ }
\newcommand{\cf}{cf.\ }
\newcommand{\resp}{respectively\ }
\newcommand{\Subsection}{\subsection}
\providecommand{\nobreakdash}{\nobreak\hbox{-}}
\setlength{\emergencystretch}{2em}
\multlinegap0pt
\usepackage{bm}
\usepackage{bbm}
\usepackage{dsfont}
\usepackage{xspace}
\usepackage{cancel}
\usepackage{mathtools}
\usepackage{amsthm}
\makeatletter
\@ifundefined{theorem}{\newtheorem{theorem}{Theorem}}{}
\makeatother
\title[A Denjoy-Wolff theorem for bounded symmetric domains]{A Denjoy-Wolff theorem for bounded symmetric domains}
\author{Cho-Ho Chu}
\date{Source: arXiv:2508.05767v1 (CC BY 4.0)}
\begin{document}
\maketitle

\noindent\emph{Source and licence.} A Denjoy-Wolff theorem for bounded symmetric domains. Version arXiv:2508.05767v1 (CC BY 4.0), distributed under the Creative Commons Attribution 4.0 International licence, as recorded in the arXiv licence metadata for this exact version and shown on the abstract page https://arxiv.org/abs/2508.05767v1. Full rights notice: licenses/arxiv-2508.05767-CC-BY-4.0.txt.\par\smallskip

\noindent\textit{Editorial scope.} Counted unit: the Theorem whose source label is \texttt{cr1} in the article (source0.tex lines 1551--1562, source label \texttt{cr1}), delivered together with the article's own complete original proof of it (source0.tex lines 1564--1608, one \texttt{proof} environment of 45 lines). No mathematics is written, repaired or re-derived: statement and proof are the article's own text line for line. Required context retained uncounted: sigi (lines 1065--1068); xi1 (lines 1073--1075); horo (lines 1234--1247); e=xi (lines 1419--1424); clos (lines 1530--1532). Every definition, display and label of those lines is retained unchanged. Omitted from this package and not redistributed: 1--1064, 1069--1072, 1076--1233, 1248--1418, 1425--1529, 1533--1550, 1563--1563, 1609--2388. Nothing inside a retained excerpt is omitted or altered: every retained body is delivered exactly as the article's lines are written, so no omission receipt of any kind accompanies this package. Citations: the retained excerpts carry no citation command at all, so no bibliography entry of the article is transcribed and no reference is redistributed. Reference adaptation: none was needed and none was made; every reference inside the delivered unit resolves inside it, so no unresolved reference or citation is produced. Printed slips kept literal: no printed word, symbol, hypothesis or number of the counted statement or of its proof was altered. Layout and class: the article's own class is \texttt{article} with packages including amsmath, amsthm, amssymb, times, enumerate, color; the wrapper is the lane's pinned \texttt{amsart} editorial wrapper at 10pt on A4, which restates the theorem-like environments the retained text needs and adds the article's own macro definitions that the retained text needs (none), with extra wrapper packages (bm, bbm, dsfont, xspace, cancel, mathtools). The delivered numbering is the unit's own. Assets: no retained span contains a figure environment, a graphics-inclusion command, a PDF-page inclusion, a table or a TikZ picture; no image file is copied into this package, the package declares no asset dependency and no image file is redistributed with it. Licence: version arXiv:2508.05767v1 of this article is distributed under the Creative Commons Attribution 4.0 International licence, as recorded in the arXiv licence metadata for this exact version and shown on the abstract page https://arxiv.org/abs/2508.05767v1. A full rights notice is delivered at licenses/arxiv-2508.05767-CC-BY-4.0.txt. Characters: every retained excerpt is pure ASCII, so the delivered engine, XeLaTeX with its default Latin Modern text font, reports no missing character.
\begin{equation}\label{sigi}
\alpha_i = \lim_k \alpha_{ki}, ~ \sigma_i =\lim_{k\rightarrow \infty} \frac{1-\alpha_{k1}^2}{ 1-\alpha_{ki}^2}
~~{\rm exist, ~ and~}(e_{ki}) ~{\rm weakly~ converges~to~} e_i
\end{equation}

\begin{equation}\label{xi1}
 F_\xi (x) =\lim_{k\rightarrow \infty} \frac{1-\|z_k\|^2}{1-\|g_{-x}(z_k)\|^2} \qquad (x\in D)
 \end{equation}

\begin{theorem}\label{horo} Let $D$ be a bounded symmetric domain of finite rank $r$, realised as the open unit
ball of a JB*-triple $V$. Then the horofunction $F_\xi$  in (\ref{xi1}), defined by   a sequence $(z_k)$ in $D$ converging to a boundary point $\xi \in \partial D$ and satisfying (\ref{sigi}), can be expressed as
\begin{equation}\label{fxi}
F_\xi(x) =\left\|\sum_{1\leq i\leq j\leq p}\rho_i\rho_jB(x,x)^{-1/2}B(x,\xi)P_{ij}\right\|
\qquad (x\in D)
\end{equation}
for some $p\in \{1, \ldots, r\}$ and $\rho_1, \ldots, \rho_p \in [0,1]$,
where  $\rho_1=1$
and  $\xi$ has a spectral decomposition
$$\xi = \alpha_1 e_1 + \alpha_2 e_2 + \cdots \alpha_p e_p \qquad (1=\alpha_1 \geq \alpha_2 \geq \cdots \geq \alpha_p>0)
$$
with joint Peirce projections $P_{ij}: V \longrightarrow V$  induced
by the minimal tripotents $e_1, \ldots, e_{p}$.
\end{theorem}

 \begin{eqnarray}\label{e=xi}
H(c,s) &=&  \sum_{j=1}^q \frac{\sigma_j}{\sigma_j+s} e_j +
B\left (\sum_{j=1}^q \sqrt{\frac{\sigma_j}{\sigma_j+s}} e_j, \,\sum_{j=1}^q \sqrt{\frac{\sigma_j}{\sigma_j +s}} e_j
\right)^{1/2}(D)\nonumber\\
&=& H(\xi,s)
\end{eqnarray}

\begin{equation}\label{clos}
\overline H (\xi,s) = c_s + B_s^{1/2}(\overline D) = \{x\in \overline D : \|B_s^{-1/2}(x-c_s)\|\leq 1\}.
\end{equation}

\begin{theorem}\label{cr1}
Let $D$ be a bounded symmetric domain of rank $r< \infty$, realised as the open
unit ball of a JB*-triple $V$. Let $(z_k)$ be a sequence in $D$ norm convergent to $\xi\in \partial D$,
which defines the horoballs
$$H(\xi, s) = \{x\in D: F_\xi(x)<s\}  \qquad (s>0)$$
with horofunction
$F_\xi$.
Then the intersection $\bigcap_{s>0}\overline H (\xi,s)$
is the closure of a holomorphic boundary component of $D$. More explicitly,
$$\bigcap_{s>0}\overline H (\xi,s)= c + (V_0(c) \cap \overline D) = \overline \Gamma_c$$
where $c\in \partial D$ is the horocentre of $H(\xi,s)$  and $V_0(c)$ is the Peirce $0$-space of $c$.
\end{theorem}

\begin{proof} Let $\xi= \alpha_1 e_1 + \cdots + \alpha_p e_p$ ,  with $\alpha_p >0$
and $p\leq r$, as in Theorem \ref{horo}. By  (\ref{e=xi}) and (\ref{clos}), we have
$$\overline H(\xi,s)=\sum_{j=1}^q \frac{\sigma_j}{\sigma_j +s} e_j +
B\left (\sum_{j=1}^q \sqrt{\frac{\sigma_j}{\sigma_j+s}} e_j, \,\sum_{j=1}^q \sqrt{\frac{\sigma_j}{\sigma_j+s}} e_j
\right)^{1/2}(\overline D)$$
where $1=\sigma_1 \geq \cdots \geq \sigma_q > 0$.
Let
 \begin{equation}\label{sume}
 c= e_1 + \cdots +e_q = \lim_{s\rightarrow 0}\sum_{j=1}^q \frac{\sigma_j}{\sigma_j+s} e_j
= \lim_{s\rightarrow 0}\sum_{j=1}^q \sqrt{\frac{\sigma_j}{s+\sigma_j}}\, e_j
\end{equation}
which is the horocentre of $H(\xi,s)$.

We show $$\bigcap_{s>0}\overline H (\xi,s) = \overline \Gamma_c.$$

Let $x\in \bigcap_{s>0}\overline H (\xi,s)$. Then for $n=1,2,\ldots,$
we have $x \in \overline H(\xi,1/n)$ and
$$ x= \sum_{j=1}^q \frac{\sigma_j}{\sigma_j + 1/n} e_j +
B\left (\sum_{j=1}^q \sqrt{\frac{\sigma_j}{\sigma_j + 1/n}} e_j, \,\sum_{j=1}^q \sqrt{\frac{\sigma_j}{\sigma_j + 1/n}} e_j
\right)^{1/2}(x_n)$$
for some $x_n \in \overline D$. By relative weak compactness of $D$, there is a subsequence $(x_k)$ of $(x_n)$
weakly converging to some point $w\in \overline D$. It follows that, letting $k\rightarrow \infty$,
$$x = c\, + \, weak\mbox{-}\lim_{k\rightarrow \infty}\,
B\hspace{-.03in}\left(\sum_{j=1}^q \sqrt{\frac{\sigma_j}{\sigma_j + 1/k}} e_j, \,\sum_{j=1}^q \sqrt{\frac{\sigma_j}{\sigma_j + 1/k}} e_j
\right)^{1/2}(x_k)$$
by (\ref{sume}), where the sequence of Bergman operators is norm convergent,  with limit
 $$\lim_{k\rightarrow \infty}\,
B\hspace{-.03in}\left(\sum_{j=1}^q \sqrt{\frac{\sigma_j}{\sigma_j + 1/k}} e_j, \,\sum_{j=1}^q \sqrt{\frac{\sigma_j}{\sigma_j + 1/k}} e_j
\right)^{1/2} =B(c,c)^{1/2}=P_0(c).$$
 Hence $x= c+ P_0(c)w \in c+ (V_0(c)\cap \overline D) = \overline \Gamma_c$.

 Conversely, let $y\in \overline \Gamma_c$. Then there is a point $z\in \overline D$
 such that
$$y= c+ P_0(c)(z) = c+ B(c,c)^{1/2}(z).$$
For $n=1, 2, \ldots$, let
$$y_n =   \sum_{j=1}^q \frac{\sigma_j}{\sigma_j + 1/n} e_j  + B\hspace{-.03in}\left(\sum_{j=1}^q \sqrt{\frac{\sigma_j}{\sigma_j + 1/n}} e_j, \,\sum_{j=1}^q \sqrt{\frac{\sigma_j}{\sigma_j + 1/n}} e_j
\right)^{1/2}(z) \in \overline H (\xi,1/n). $$
Then we have $ \lim_n y_n =y$.

If $y \notin \bigcap_{s>0}\overline H (\xi,s)$, then $y \notin\overline H (\xi,s_0)$ for some $s_0 > 0$.
Hence there exists $N \in \mathbb{N}$ such that $y_n \notin\overline H (\xi,s_0)$ for all $n >N$.
Pick $n'>N$ such that $1/n' < s_0$. Then $y_{n'} \notin\overline H (\xi,s_0)$ whereas
$y_{n'} \in\overline H (\xi, 1/n') \subset \overline H(\xi,s_0)$, which is impossible. Therefore $y \in
\bigcap_{s>0}\overline H (\xi, s)$ and the proof is complete.
\end{proof}

\end{document}
